\uuid{kDGn}
\titre{Changement de variables}
\auteur{hanani ; mzari}
\organisation{exo7}
\licence{CC BY-NC-SA 4.0}
\datecreate{2020-01-20}
\niveau{L1}
\module{Analyse}
\chapitre{Calcul d'intégrales}
\sousChapitre{Primitives diverses}
\difficulte{3}
\theme{Primitives, changement de variables}
\exo7qcmid{lille-2.primitives.030}

\contenu{
\begin{question}
\qclassification{127.11}
\qtags{motcle=changement de variables}
Parmi les égalités suivantes, cocher celles qui sont vraies :
\begin{answers}
\bad{$\displaystyle \int \cos (2x)\sqrt{1+\sin (2x)}\mathrm{d}x=\frac{\cos (2x)}{\sqrt{1+\sin (2x)}}+k$, où $k\in \Rr$.}
\good{$\displaystyle \int \cos (2x)\sqrt{1+\sin (2x)}\mathrm{d}x=\frac{1}{3}\left[1+\sin (2x)\right]^{3/2}+k$, où $k\in \Rr$.}
\good{$\displaystyle \int \frac{\sin (3x)\mathrm{d}x}{2-\cos (3x)}=\frac{1}{3}\ln \left[2-\cos (3x)\right]+k$ où $k\in \Rr$.}
\bad{$\displaystyle \int \frac{\mathrm{d}x}{2-\cos (3x)}=\ln \left[2-\cos (3x)\right]+k$, $k\in \Rr$.}
\end{answers}
\begin{explanations}
Avec $u=1+\sin (2x)$, on aura $\displaystyle \mathrm{d}u=2\cos (2x)\mathrm{d}x$. D'où
$$\int \cos (2x)\sqrt{1+\sin (2x)}\mathrm{d}x=2\int \sqrt{u}\mathrm{d}u=\frac{1}{3}u^{3/2}+k,\; k\in \Rr.$$
De même, avec $u=2-\cos (3x)$, on aura $\displaystyle \mathrm{d}u=3\sin(3x)\mathrm{d}x$. D'où
$$\int \frac{\sin (3x)\mathrm{d}x}{2-\cos (3x)}=\frac{1}{3}\int \frac{\mathrm{d}u}{u}=\frac{1}{3}\ln |u|+k,\; k\in \Rr.$$
\end{explanations}
\end{question}
}
